% Options for packages loaded elsewhere
\PassOptionsToPackage{unicode}{hyperref}
\PassOptionsToPackage{hyphens}{url}
%
\documentclass[
]{ctexart}
\usepackage{amsmath,amssymb}
\usepackage{iftex}
\ifPDFTeX
  \usepackage[T1]{fontenc}
  \usepackage[utf8]{inputenc}
  \usepackage{textcomp} % provide euro and other symbols
\else % if luatex or xetex
  \usepackage{unicode-math} % this also loads fontspec
  \defaultfontfeatures{Scale=MatchLowercase}
  \defaultfontfeatures[\rmfamily]{Ligatures=TeX,Scale=1}
\fi
\usepackage{lmodern}
\ifPDFTeX\else
  % xetex/luatex font selection
\fi
% Use upquote if available, for straight quotes in verbatim environments
\IfFileExists{upquote.sty}{\usepackage{upquote}}{}
\IfFileExists{microtype.sty}{% use microtype if available
  \usepackage[]{microtype}
  \UseMicrotypeSet[protrusion]{basicmath} % disable protrusion for tt fonts
}{}
\makeatletter
\@ifundefined{KOMAClassName}{% if non-KOMA class
  \IfFileExists{parskip.sty}{%
    \usepackage{parskip}
  }{% else
    \setlength{\parindent}{0pt}
    \setlength{\parskip}{6pt plus 2pt minus 1pt}}
}{% if KOMA class
  \KOMAoptions{parskip=half}}
\makeatother
\usepackage{xcolor}
\usepackage{color}
\usepackage{fancyvrb}
\newcommand{\VerbBar}{|}
\newcommand{\VERB}{\Verb[commandchars=\\\{\}]}
\DefineVerbatimEnvironment{Highlighting}{Verbatim}{commandchars=\\\{\}}
% Add ',fontsize=\small' for more characters per line
\usepackage{framed}
\definecolor{shadecolor}{RGB}{248,248,248}
\newenvironment{Shaded}{\begin{snugshade}}{\end{snugshade}}
\newcommand{\AlertTok}[1]{\textcolor[rgb]{0.94,0.16,0.16}{#1}}
\newcommand{\AnnotationTok}[1]{\textcolor[rgb]{0.56,0.35,0.01}{\textbf{\textit{#1}}}}
\newcommand{\AttributeTok}[1]{\textcolor[rgb]{0.13,0.29,0.53}{#1}}
\newcommand{\BaseNTok}[1]{\textcolor[rgb]{0.00,0.00,0.81}{#1}}
\newcommand{\BuiltInTok}[1]{#1}
\newcommand{\CharTok}[1]{\textcolor[rgb]{0.31,0.60,0.02}{#1}}
\newcommand{\CommentTok}[1]{\textcolor[rgb]{0.56,0.35,0.01}{\textit{#1}}}
\newcommand{\CommentVarTok}[1]{\textcolor[rgb]{0.56,0.35,0.01}{\textbf{\textit{#1}}}}
\newcommand{\ConstantTok}[1]{\textcolor[rgb]{0.56,0.35,0.01}{#1}}
\newcommand{\ControlFlowTok}[1]{\textcolor[rgb]{0.13,0.29,0.53}{\textbf{#1}}}
\newcommand{\DataTypeTok}[1]{\textcolor[rgb]{0.13,0.29,0.53}{#1}}
\newcommand{\DecValTok}[1]{\textcolor[rgb]{0.00,0.00,0.81}{#1}}
\newcommand{\DocumentationTok}[1]{\textcolor[rgb]{0.56,0.35,0.01}{\textbf{\textit{#1}}}}
\newcommand{\ErrorTok}[1]{\textcolor[rgb]{0.64,0.00,0.00}{\textbf{#1}}}
\newcommand{\ExtensionTok}[1]{#1}
\newcommand{\FloatTok}[1]{\textcolor[rgb]{0.00,0.00,0.81}{#1}}
\newcommand{\FunctionTok}[1]{\textcolor[rgb]{0.13,0.29,0.53}{\textbf{#1}}}
\newcommand{\ImportTok}[1]{#1}
\newcommand{\InformationTok}[1]{\textcolor[rgb]{0.56,0.35,0.01}{\textbf{\textit{#1}}}}
\newcommand{\KeywordTok}[1]{\textcolor[rgb]{0.13,0.29,0.53}{\textbf{#1}}}
\newcommand{\NormalTok}[1]{#1}
\newcommand{\OperatorTok}[1]{\textcolor[rgb]{0.81,0.36,0.00}{\textbf{#1}}}
\newcommand{\OtherTok}[1]{\textcolor[rgb]{0.56,0.35,0.01}{#1}}
\newcommand{\PreprocessorTok}[1]{\textcolor[rgb]{0.56,0.35,0.01}{\textit{#1}}}
\newcommand{\RegionMarkerTok}[1]{#1}
\newcommand{\SpecialCharTok}[1]{\textcolor[rgb]{0.81,0.36,0.00}{\textbf{#1}}}
\newcommand{\SpecialStringTok}[1]{\textcolor[rgb]{0.31,0.60,0.02}{#1}}
\newcommand{\StringTok}[1]{\textcolor[rgb]{0.31,0.60,0.02}{#1}}
\newcommand{\VariableTok}[1]{\textcolor[rgb]{0.00,0.00,0.00}{#1}}
\newcommand{\VerbatimStringTok}[1]{\textcolor[rgb]{0.31,0.60,0.02}{#1}}
\newcommand{\WarningTok}[1]{\textcolor[rgb]{0.56,0.35,0.01}{\textbf{\textit{#1}}}}
\usepackage{graphicx}
\makeatletter
\def\maxwidth{\ifdim\Gin@nat@width>\linewidth\linewidth\else\Gin@nat@width\fi}
\def\maxheight{\ifdim\Gin@nat@height>\textheight\textheight\else\Gin@nat@height\fi}
\makeatother
% Scale images if necessary, so that they will not overflow the page
% margins by default, and it is still possible to overwrite the defaults
% using explicit options in \includegraphics[width, height, ...]{}
\setkeys{Gin}{width=\maxwidth,height=\maxheight,keepaspectratio}
% Set default figure placement to htbp
\makeatletter
\def\fps@figure{htbp}
\makeatother
\setlength{\emergencystretch}{3em} % prevent overfull lines
\providecommand{\tightlist}{%
  \setlength{\itemsep}{0pt}\setlength{\parskip}{0pt}}
\setcounter{secnumdepth}{5}
\renewcommand{\and}{\\}
\ifLuaTeX
  \usepackage{selnolig}  % disable illegal ligatures
\fi
\IfFileExists{bookmark.sty}{\usepackage{bookmark}}{\usepackage{hyperref}}
\IfFileExists{xurl.sty}{\usepackage{xurl}}{} % add URL line breaks if available
\urlstyle{same}
\hypersetup{
  pdftitle={Correlation \& Distance},
  pdfauthor={Rui Zhou},
  hidelinks,
  pdfcreator={LaTeX via pandoc}}

\title{Correlation \& Distance}
\author{Rui Zhou}
\date{2026-05-26}

\begin{document}
\maketitle

{
\setcounter{tocdepth}{2}
\tableofcontents
}
\begin{Shaded}
\begin{Highlighting}[]
\FunctionTok{library}\NormalTok{(bruceR)}
\FunctionTok{library}\NormalTok{(rstatix)}
\FunctionTok{library}\NormalTok{(tidyverse)}
\FunctionTok{library}\NormalTok{(stargazer)}
\end{Highlighting}
\end{Shaded}

\hypertarget{correlation}{%
\section{Correlation}\label{correlation}}

\hypertarget{bivariate-correlation}{%
\subsection{Bivariate Correlation}\label{bivariate-correlation}}

Here we're to discuss more about bivariate correlation when taking
groups into account. As you can see, the pooled data may interact and
neutralize the effect or relationship of two variables. However, after
we inspect that correlation under each group, surprising outcomes may
spring up!

Take dataset ``car\_sales.sav'' as an example, which you have been quite
similar to ``mtcars''. We're interested in whether or not people value
fuel-efficient cars more.

\begin{Shaded}
\begin{Highlighting}[]
\NormalTok{carsales }\OtherTok{=} \FunctionTok{import}\NormalTok{(}\StringTok{"car\_sales.sav"}\NormalTok{)}
\end{Highlighting}
\end{Shaded}

\begin{verbatim}
## ✔ Successfully imported: 157 obs. of 26 variables
\end{verbatim}

\begin{Shaded}
\begin{Highlighting}[]
\FunctionTok{attach}\NormalTok{(carsales)}
\end{Highlighting}
\end{Shaded}

\begin{Shaded}
\begin{Highlighting}[]
\FunctionTok{ggplot}\NormalTok{(}\AttributeTok{data =}\NormalTok{ carsales, }\AttributeTok{mapping =} \FunctionTok{aes}\NormalTok{(}\AttributeTok{x =}\NormalTok{ mpg,}\AttributeTok{y =}\NormalTok{ sales))}\SpecialCharTok{+}
  \FunctionTok{geom\_point}\NormalTok{()}
\end{Highlighting}
\end{Shaded}

\begin{figure}

{\centering \includegraphics[width=0.7\linewidth]{r-correlation-distance_files/figure-latex/unnamed-chunk-3-1} 

}

\caption{Scatter Plot of mpg and sales with ALL Observations}\label{fig:unnamed-chunk-3}
\end{figure}

At first glance, there are apparently some outliers. Moreover, the
scatter points seem to cluster in the bottom-left, possibly a skewed
distribution instead of normal! Nevertheless, those don't bother us to
check the overall correlation first.

\begin{Shaded}
\begin{Highlighting}[]
\FunctionTok{cor}\NormalTok{(sales,mpg,}\AttributeTok{use =} \StringTok{\textquotesingle{}c\textquotesingle{}}\NormalTok{)}
\end{Highlighting}
\end{Shaded}

\begin{verbatim}
## [1] -0.01668058
\end{verbatim}

The correlation is not high in magnitude. Here, use ``\texttt{cor()}''
for simplicity. Note that for data with missing values, specify
parameter
``\texttt{use\ =\ \textquotesingle{}complete.obs\textquotesingle{}}''
(or ``\texttt{use\ =\ \textquotesingle{}c\textquotesingle{}}'' as
abbreviation) to compute only with complete cases.

Admittedly, ``\texttt{cor.test()}'' will also return the correlation
coefficient, along with outcomes with hypothesis testing. However, the
latter makes it to be complex unnecessarily! (Good thing is that,
``\texttt{cor.test()}'' will automatically wipe out cases with missing
values.)

\begin{Shaded}
\begin{Highlighting}[]
\FunctionTok{cor.test}\NormalTok{(mpg,sales)}
\end{Highlighting}
\end{Shaded}

Then, we're to deal with outliers and skewness, then to recompute
correlation.

Good news is that, \(\log\) transformation helps when it comes to skewed
distribution, and outliers will get contained under \(\log\)
transformation at the same time.

\begin{Shaded}
\begin{Highlighting}[]
\FunctionTok{cor}\NormalTok{(}\FunctionTok{log}\NormalTok{(mpg),}\FunctionTok{log}\NormalTok{(sales),}\AttributeTok{use =} \StringTok{\textquotesingle{}c\textquotesingle{}}\NormalTok{)}
\end{Highlighting}
\end{Shaded}

\begin{verbatim}
## [1] 0.1141838
\end{verbatim}

The correlation coefficient is improved in magnitude, but not
significant after all.

Then, here comes what's the most important point. When faced with a
brand-new data, before conducting any computation and hypothesis
testing, getting an overall picture of the data is the most important
thing to do!

After observation, it's likely that cars whose model are ``Metro'' or
``F-Series'' are special compared to others. Accordingly, we're to group
the data and then recompute correlation.

\begin{Shaded}
\begin{Highlighting}[]
\NormalTok{sales\_some }\OtherTok{=}\NormalTok{ sales[model}\SpecialCharTok{!=}\StringTok{"Metro"}\SpecialCharTok{\&}\NormalTok{model}\SpecialCharTok{!=}\StringTok{"F{-}Series"}\NormalTok{]}
\NormalTok{mpg\_some }\OtherTok{=}\NormalTok{ mpg[model}\SpecialCharTok{!=}\StringTok{"Metro"}\SpecialCharTok{\&}\NormalTok{model}\SpecialCharTok{!=}\StringTok{"F{-}Series"}\NormalTok{]}
\NormalTok{type\_some }\OtherTok{=} \FunctionTok{as.factor}\NormalTok{(type[model}\SpecialCharTok{!=}\StringTok{"Metro"}\SpecialCharTok{\&}\NormalTok{model}\SpecialCharTok{!=}\StringTok{"F{-}Series"}\NormalTok{])}

\NormalTok{df }\OtherTok{=} \FunctionTok{data.frame}\NormalTok{(sales\_some,mpg\_some,type\_some)}
\NormalTok{df }\SpecialCharTok{\%\textgreater{}\%}
  \FunctionTok{ggplot}\NormalTok{(}\FunctionTok{aes}\NormalTok{(mpg\_some, sales\_some,}\AttributeTok{color =}\NormalTok{ type\_some)) }\SpecialCharTok{+}
  \FunctionTok{geom\_point}\NormalTok{(}\AttributeTok{alpha=}\FloatTok{0.8}\NormalTok{, }\AttributeTok{size=}\FloatTok{0.8}\NormalTok{) }\SpecialCharTok{+}
  \FunctionTok{geom\_smooth}\NormalTok{(}\AttributeTok{method=}\StringTok{"lm"}\NormalTok{,}\AttributeTok{se=}\ConstantTok{FALSE}\NormalTok{) }\SpecialCharTok{+}
  \FunctionTok{ylab}\NormalTok{(}\StringTok{\textquotesingle{}Sales(k)\textquotesingle{}}\NormalTok{) }\SpecialCharTok{+}
  \FunctionTok{xlab}\NormalTok{(}\StringTok{\textquotesingle{}MPG\textquotesingle{}}\NormalTok{)}
\end{Highlighting}
\end{Shaded}

\includegraphics{r-correlation-distance_files/figure-latex/unnamed-chunk-7-1.pdf}

Use graphs to think! The categorized data plays an important role here,
and unfolds the underlying correlation. Based on the prior observation,
then check correlation quantitatively.

\begin{Shaded}
\begin{Highlighting}[]
\FunctionTok{cor.test}\NormalTok{(sales\_some[type\_some}\SpecialCharTok{==}\StringTok{"1"}\NormalTok{],mpg\_some[type\_some}\SpecialCharTok{==}\StringTok{"1"}\NormalTok{])}
\end{Highlighting}
\end{Shaded}

\begin{verbatim}
## 
##  Pearson's product-moment correlation
## 
## data:  sales_some[type_some == "1"] and mpg_some[type_some == "1"]
## t = 0.070377, df = 38, p-value = 0.9443
## alternative hypothesis: true correlation is not equal to 0
## 95 percent confidence interval:
##  -0.3011645  0.3217809
## sample estimates:
##        cor 
## 0.01141584
\end{verbatim}

\begin{Shaded}
\begin{Highlighting}[]
\FunctionTok{cor.test}\NormalTok{(sales\_some[type\_some}\SpecialCharTok{==}\StringTok{"0"}\NormalTok{],mpg\_some[type\_some}\SpecialCharTok{==}\StringTok{"0"}\NormalTok{])}
\end{Highlighting}
\end{Shaded}

\begin{verbatim}
## 
##  Pearson's product-moment correlation
## 
## data:  sales_some[type_some == "0"] and mpg_some[type_some == "0"]
## t = 2.8695, df = 112, p-value = 0.004915
## alternative hypothesis: true correlation is not equal to 0
## 95 percent confidence interval:
##  0.0817139 0.4251484
## sample estimates:
##       cor 
## 0.2616958
\end{verbatim}

Finally, don't forget to

\begin{Shaded}
\begin{Highlighting}[]
\FunctionTok{detach}\NormalTok{(carsales)}
\end{Highlighting}
\end{Shaded}

\hypertarget{partial-part-correlation}{%
\subsection{Partial \& Part
Correlation}\label{partial-part-correlation}}

\hypertarget{background-knowledge}{%
\subsubsection{Background Knowledge}\label{background-knowledge}}

Above all, I'd like to clarify one concept that we've been familiar with
but without mentioning this concept. Zero-order correlation refers to
the correlation between \(X\) and \(Y\) without controlling other
covariates. Traditionally, Pearson and Spearman correlations are typical
of zero-order correlation.

We've discussed partial and part (or, semi-partial) correlation in
multilinear regression. Take a step further, we can consider those under
the most general cases, i.e., without the setting of MLR but some
variables.

Suppose we have three variables, called \(X,Y\) and \(Z\), where \(X\)
is the predictor, \(Y\) is the dependent variable, and \(Z\) is the
covariate. Noteworthy that the assumption of the roles of \(X,Y,Z\)
don't matter much, only for accommodating the setting.

\begin{center}\includegraphics[width=0.6\linewidth]{partial & part correlation} \end{center}

Note that, in the general case, \(Y\) is not represented as the whole
area in the picture, but just an equal role as a circle.

If the correlation of \(X\) and \(Y\) is of our great interest, \[
\begin{aligned}
r_x&=\frac {a+c}{a+b+c+d}\\
pr_x&=\frac{a}{a+d}\\
spr_x&=\frac{a}{a+b+c+d}
\end{aligned}
\]

\begin{itemize}
\item
  Partial: Holding \(Z\) from BOTH \(X\) and \(Y\), i.e., removing the
  effect of \(Z\).

  Good for knowing the \textbf{variance uncounted by other control
  variable(s)}.
\item
  Semi-Partial: Holding \(Z\) from \(X\) only.

  Good for knowing the \textbf{unique contribution of \(X\)}.
\item
  Zero-Order: No extra constraint.
\end{itemize}

One more enlightening way to think about partial correlation is through
multilinear regression.

\begin{center}\includegraphics[width=0.4\linewidth]{MLR partial corr} \end{center}

If we're interested in the partial correlation between \(A\) and \(B\),
with \(C\) as the covariate. Then, what we need to do is partial out the
effect of \(C\) on BOTH \(A\) and \(B\), and this goal can be
accomplished through step-by-step linear regression.

Specifically, first we conduct a simple linear regression of \(A\) on
\(C\), and the residual part is that of \(A\) which doesn't intersect
with \(C\). Similarly, then followed by a SLR of \(B\) on \(C\), same
with the meaning of residual part. For the two residual parts, they
share it in common that they're clean of \(C\). Lastly, compute the
correlation between \(A\) and \(B\) directly, and that's what we want!

Even with more covariates, conduct \(A\) and \(B\) on all the covariates
and then conpute the correlation with their residual parts. Done!

\hypertarget{examples}{%
\subsubsection{Examples}\label{examples}}

We first man-make a data that vividly shows the differences between
partial and zero-order correlation.

\begin{Shaded}
\begin{Highlighting}[]
\NormalTok{DV }\OtherTok{=} \FunctionTok{c}\NormalTok{(}\DecValTok{1}\SpecialCharTok{:}\DecValTok{8}\NormalTok{, }\DecValTok{8}\SpecialCharTok{:}\DecValTok{11}\NormalTok{)}
\NormalTok{IV }\OtherTok{=} \FunctionTok{c}\NormalTok{(}\DecValTok{4}\SpecialCharTok{:}\DecValTok{1}\NormalTok{,}\DecValTok{8}\SpecialCharTok{:}\DecValTok{5}\NormalTok{,}\DecValTok{12}\SpecialCharTok{:}\DecValTok{9}\NormalTok{)}
\NormalTok{covariate }\OtherTok{=} \FunctionTok{c}\NormalTok{(}\FunctionTok{rep}\NormalTok{(}\DecValTok{0}\NormalTok{,}\DecValTok{4}\NormalTok{), }\FunctionTok{rep}\NormalTok{(}\DecValTok{4}\NormalTok{,}\DecValTok{4}\NormalTok{), }\FunctionTok{rep}\NormalTok{(}\DecValTok{7}\NormalTok{,}\DecValTok{4}\NormalTok{))}
\NormalTok{studiedData }\OtherTok{=} \FunctionTok{data.frame}\NormalTok{(DV,IV,covariate)}

\NormalTok{studiedData }\SpecialCharTok{|\textgreater{}} \FunctionTok{ggplot}\NormalTok{(}\FunctionTok{aes}\NormalTok{(}\AttributeTok{x =}\NormalTok{ IV,}\AttributeTok{y =}\NormalTok{ DV))}\SpecialCharTok{+}
  \FunctionTok{geom\_point}\NormalTok{()}
\end{Highlighting}
\end{Shaded}

\begin{figure}

{\centering \includegraphics[width=0.7\linewidth]{r-correlation-distance_files/figure-latex/unnamed-chunk-12-1} 

}

\caption{Scatter Plot of X and Y}\label{fig:unnamed-chunk-12}
\end{figure}

Clearly, considering zero-order correlation, \(X\) is strictly
negatively-correlated with \(Y\), the covariate is strictly
positively-correlated with \(Y\). However, \(X\) shows partially
positive correlation with \(Y\) overall.

\begin{Shaded}
\begin{Highlighting}[]
\FunctionTok{cor}\NormalTok{(IV,DV)}
\end{Highlighting}
\end{Shaded}

\begin{verbatim}
## [1] 0.7608292
\end{verbatim}

If you're to study the partial correlation, things will be different.

\begin{Shaded}
\begin{Highlighting}[]
\FunctionTok{library}\NormalTok{(ppcor)}
\FunctionTok{pcor}\NormalTok{(studiedData)}
\end{Highlighting}
\end{Shaded}

\begin{table}[!htbp] \centering 
  \caption{Partial Correlation between Variables} 
  \label{} 
\begin{tabular}{@{\extracolsep{5pt}} cccc} 
\\[-1.8ex]\hline 
\hline \\[-1.8ex] 
 & DV & IV & covariate \\ 
\hline \\[-1.8ex] 
DV & $1$ & $$-$0.972$ & $0.991$ \\ 
IV & $$-$0.972$ & $1$ & $0.993$ \\ 
covariate & $0.991$ & $0.993$ & $1$ \\ 
\hline \\[-1.8ex] 
\end{tabular} 
\end{table}

One more example. A popular radio talk show host has just received the
latest government study on public health care funding and has uncovered
a startling fact: As health care funding increases, disease rates also
increase! Cities that spend more actually seem to be worse off than
cities that spend less. Then, shall we just cut the funding and save
people's health?

\begin{Shaded}
\begin{Highlighting}[]
\CommentTok{\# example: a covariate can "make" a correlation that should not exist}
\NormalTok{currentData }\OtherTok{\textless{}{-}} \FunctionTok{import}\NormalTok{(}\StringTok{"health\_funding.sav"}\NormalTok{)}
\end{Highlighting}
\end{Shaded}

\begin{verbatim}
## ✔ Successfully imported: 50 obs. of 4 variables
\end{verbatim}

\begin{Shaded}
\begin{Highlighting}[]
\FunctionTok{attach}\NormalTok{(currentData)}
\NormalTok{studiedData }\OtherTok{=} \FunctionTok{data.frame}\NormalTok{(disease,funding,visits)}
\end{Highlighting}
\end{Shaded}

First of all, check overall zero-order correlation.

\begin{Shaded}
\begin{Highlighting}[]
\FunctionTok{cor}\NormalTok{(studiedData)}
\end{Highlighting}
\end{Shaded}

\begin{table}[!htbp] \centering 
  \caption{Zero-Order Correlation} 
  \label{} 
\begin{tabular}{@{\extracolsep{5pt}} cccc} 
\\[-1.8ex]\hline 
\hline \\[-1.8ex] 
 & disease & funding & visits \\ 
\hline \\[-1.8ex] 
disease & $1$ & $0.737$ & $0.762$ \\ 
funding & $0.737$ & $1$ & $0.964$ \\ 
visits & $0.762$ & $0.964$ & $1$ \\ 
\hline \\[-1.8ex] 
\end{tabular} 
\end{table}

It seems astonishing that ``\texttt{funding}'' is greatly correlated
with ``\texttt{disease}''! May it offer the plausible reason to cut the
funding? Remember, we still have the covariate ``\texttt{visits}''. One
alternative explanation might be that, offered with more funding, people
may be more likely to pay visits to health care, and thus more diseases
are detected, which should have been detected but not due to wealth
state.

\begin{Shaded}
\begin{Highlighting}[]
\FunctionTok{pcor}\NormalTok{(studiedData)}\SpecialCharTok{$}\NormalTok{estimate}
\end{Highlighting}
\end{Shaded}

\begin{table}[!htbp] \centering 
  \caption{Partial Correlation} 
  \label{} 
\begin{tabular}{@{\extracolsep{5pt}} cccc} 
\\[-1.8ex]\hline 
\hline \\[-1.8ex] 
 & disease & funding & visits \\ 
\hline \\[-1.8ex] 
disease & $1$ & $0.013$ & $0.286$ \\ 
funding & $0.013$ & $1$ & $0.920$ \\ 
visits & $0.286$ & $0.920$ & $1$ \\ 
\hline \\[-1.8ex] 
\end{tabular} 
\end{table}

It's astonishing that, considering the partial correlation,
``\texttt{funding}'' almost has no correlation with ``\texttt{disease}''
in essence!

Additionally, we can consider the part correlation.

\begin{Shaded}
\begin{Highlighting}[]
\FunctionTok{spcor}\NormalTok{(studiedData)}\SpecialCharTok{$}\NormalTok{estimate}
\end{Highlighting}
\end{Shaded}

\begin{table}[!htbp] \centering 
  \caption{Part Correlation} 
  \label{} 
\begin{tabular}{@{\extracolsep{5pt}} cccc} 
\\[-1.8ex]\hline 
\hline \\[-1.8ex] 
 & disease & funding & visits \\ 
\hline \\[-1.8ex] 
disease & $1$ & $0.009$ & $0.193$ \\ 
funding & $0.004$ & $1$ & $0.622$ \\ 
visits & $0.076$ & $0.596$ & $1$ \\ 
\hline \\[-1.8ex] 
\end{tabular} 
\end{table}

And finally,

\begin{Shaded}
\begin{Highlighting}[]
\FunctionTok{detach}\NormalTok{(currentData)}
\end{Highlighting}
\end{Shaded}

\hypertarget{distance}{%
\section{Distance}\label{distance}}

We need to distinguish between two kinds of distance. One is distance
between cases (or observations), the other is distance between
variables.

There are traditionally three ways to measure distance, either from the
perspective of dissimilarity or similarity.

\begin{itemize}
\tightlist
\item
  Dissimilarity

  \begin{itemize}
  \tightlist
  \item
    Euclidean distance
  \end{itemize}

  \[d(\vec p,\vec q)=\sqrt{\sum_{i=1}^n(p_i-q_i)^2}\]

  \begin{itemize}
  \tightlist
  \item
    Chebychev distance
  \end{itemize}

  A vector space where the distance between two vectors is the greatest
  of their differences along any coordinate dimension.
  \[D(\vec x,\vec y)=\max_i|x_i-y_i|\]
\item
  Similarity

  \begin{itemize}
  \tightlist
  \item
    Cosine similarity

    \begin{itemize}
    \tightlist
    \item
      A measure of similarity between two non-zero vectors of an inner
      product space that measures the cosine of the angle between them.
    \item
      It is thus a judgment of orientation and not magnitude.
    \end{itemize}
  \end{itemize}
\end{itemize}

\begin{Shaded}
\begin{Highlighting}[]
\NormalTok{currentData }\OtherTok{=} \FunctionTok{import}\NormalTok{(}\StringTok{"data10{-}04.sav"}\NormalTok{, }\AttributeTok{sheet =} \DecValTok{1}\NormalTok{)}
\end{Highlighting}
\end{Shaded}

\begin{verbatim}
## ✔ Successfully imported: 12 obs. of 6 variables
\end{verbatim}

\begin{Shaded}
\begin{Highlighting}[]
\FunctionTok{attach}\NormalTok{(currentData)}
\CommentTok{\# only use the data that should be considered or analyzed}
\NormalTok{studiedData }\OtherTok{=} \FunctionTok{data.frame}\NormalTok{(hgrow,temp,rain,hsun,humi)}
\end{Highlighting}
\end{Shaded}

\hypertarget{correlation-1}{%
\subsection{Correlation}\label{correlation-1}}

Note that correlation itself is a measure of distance. However,
correlation only applies for \textbf{distance between variables}, not
for distance between cases. (If you insist to test the correlation
between cases, easy, just transpose your data!)

\begin{Shaded}
\begin{Highlighting}[]
\FunctionTok{cor}\NormalTok{(studiedData)}
\end{Highlighting}
\end{Shaded}

\begin{table}[!htbp] \centering 
  \caption{Correlation between Variables} 
  \label{} 
\begin{tabular}{@{\extracolsep{5pt}} cccccc} 
\\[-1.8ex]\hline 
\hline \\[-1.8ex] 
 & hgrow & temp & rain & hsun & humi \\ 
\hline \\[-1.8ex] 
hgrow & $1$ & $0.983$ & $0.709$ & $0.704$ & $0.374$ \\ 
temp & $0.983$ & $1$ & $0.715$ & $0.690$ & $0.292$ \\ 
rain & $0.709$ & $0.715$ & $1$ & $0.702$ & $0.384$ \\ 
hsun & $0.704$ & $0.690$ & $0.702$ & $1$ & $$-$0.051$ \\ 
humi & $0.374$ & $0.292$ & $0.384$ & $$-$0.051$ & $1$ \\ 
\hline \\[-1.8ex] 
\end{tabular} 
\end{table}

\hypertarget{euclidean-distance}{%
\subsection{Euclidean Distance}\label{euclidean-distance}}

Then, we're interested in Euclidean distance \textbf{between cases} by
default. Use ``\texttt{dist(x,\ method\ =\ "euclidean")}''. Mind that
the distance makes sense only if the data is scaled!

\begin{Shaded}
\begin{Highlighting}[]
\CommentTok{\# studiedData \%\textgreater{}\% dist()}
\NormalTok{studiedData  }\SpecialCharTok{\%\textgreater{}\%} 
  \FunctionTok{scale}\NormalTok{() }\SpecialCharTok{\%\textgreater{}\%} \FunctionTok{dist}\NormalTok{() }
\end{Highlighting}
\end{Shaded}

\begin{table}[!htbp] \centering 
  \caption{Euclidean Distance between Cases} 
  \label{} 
\begin{tabular}{@{\extracolsep{5pt}} ccccccccccccc} 
\\[-1.8ex]\hline 
\hline \\[-1.8ex] 
 & 1 & 2 & 3 & 4 & 5 & 6 & 7 & 8 & 9 & 10 & 11 & 12 \\ 
\hline \\[-1.8ex] 
1 & $0$ & $0.925$ & $2.142$ & $3.692$ & $4.790$ & $3.904$ & $4.078$ & $5.793$ & $3.764$ & $2.536$ & $1.782$ & $0.591$ \\ 
2 & $0.925$ & $0$ & $1.288$ & $2.856$ & $4.179$ & $3.338$ & $3.756$ & $5.505$ & $3.576$ & $2.232$ & $1.433$ & $0.923$ \\ 
3 & $2.142$ & $1.288$ & $0$ & $2.149$ & $3.774$ & $3.151$ & $3.962$ & $5.710$ & $3.870$ & $2.675$ & $2.054$ & $2.192$ \\ 
4 & $3.692$ & $2.856$ & $2.149$ & $0$ & $2.558$ & $2.001$ & $3.108$ & $4.535$ & $3.518$ & $2.627$ & $2.509$ & $3.447$ \\ 
5 & $4.790$ & $4.179$ & $3.774$ & $2.558$ & $0$ & $2.307$ & $2.871$ & $2.895$ & $2.724$ & $3.151$ & $3.662$ & $4.501$ \\ 
6 & $3.904$ & $3.338$ & $3.151$ & $2.001$ & $2.307$ & $0$ & $1.265$ & $3.338$ & $1.937$ & $1.636$ & $2.236$ & $3.529$ \\ 
7 & $4.078$ & $3.756$ & $3.962$ & $3.108$ & $2.871$ & $1.265$ & $0$ & $2.696$ & $1.292$ & $1.566$ & $2.448$ & $3.629$ \\ 
8 & $5.793$ & $5.505$ & $5.710$ & $4.535$ & $2.895$ & $3.338$ & $2.696$ & $0$ & $2.761$ & $3.721$ & $4.519$ & $5.333$ \\ 
9 & $3.764$ & $3.576$ & $3.870$ & $3.518$ & $2.724$ & $1.937$ & $1.292$ & $2.761$ & $0$ & $1.654$ & $2.541$ & $3.402$ \\ 
10 & $2.536$ & $2.232$ & $2.675$ & $2.627$ & $3.151$ & $1.636$ & $1.566$ & $3.721$ & $1.654$ & $0$ & $0.967$ & $2.093$ \\ 
11 & $1.782$ & $1.433$ & $2.054$ & $2.509$ & $3.662$ & $2.236$ & $2.448$ & $4.519$ & $2.541$ & $0.967$ & $0$ & $1.374$ \\ 
12 & $0.591$ & $0.923$ & $2.192$ & $3.447$ & $4.501$ & $3.529$ & $3.629$ & $5.333$ & $3.402$ & $2.093$ & $1.374$ & $0$ \\ 
\hline \\[-1.8ex] 
\end{tabular} 
\end{table}

If the Euclidean distance between \textbf{variables} is of interest,
then \emph{transpose} your data.

\begin{Shaded}
\begin{Highlighting}[]
\NormalTok{dist }\OtherTok{=}\NormalTok{ studiedData  }\SpecialCharTok{\%\textgreater{}\%} 
  \FunctionTok{scale}\NormalTok{() }\SpecialCharTok{\%\textgreater{}\%} \FunctionTok{t}\NormalTok{() }\SpecialCharTok{\%\textgreater{}\%} \FunctionTok{dist}\NormalTok{()}
\end{Highlighting}
\end{Shaded}

\begin{table}[!htbp] \centering 
  \caption{Euclidean Distance between Variables} 
  \label{} 
\begin{tabular}{@{\extracolsep{5pt}} cccccc} 
\\[-1.8ex]\hline 
\hline \\[-1.8ex] 
 & hgrow & temp & rain & hsun & humi \\ 
\hline \\[-1.8ex] 
hgrow & $0$ & $0.605$ & $2.529$ & $2.550$ & $3.712$ \\ 
temp & $0.605$ & $0$ & $2.505$ & $2.609$ & $3.947$ \\ 
rain & $2.529$ & $2.505$ & $0$ & $2.561$ & $3.680$ \\ 
hsun & $2.550$ & $2.609$ & $2.561$ & $0$ & $4.808$ \\ 
humi & $3.712$ & $3.947$ & $3.680$ & $4.808$ & $0$ \\ 
\hline \\[-1.8ex] 
\end{tabular} 
\end{table}

It's noteworthy that, the Euclidean distance result returned by
``\texttt{dist()}'' is a unique type of ``\texttt{dist}'', which is
incompatible with some other types. Luckily, use
``\texttt{as.matrix()}'' to transform that into a matrix (with the
diagonal and upper right triangle filled automatically).

\hypertarget{cosine-similarity}{%
\subsection{Cosine Similarity}\label{cosine-similarity}}

Also, we can see the cosine similarity \textbf{between cases}.

Remember, if you're to study the distance between cases, transpose your
matrix of data; if you're to study the distance between variables, throw
the matrix of data into ``\texttt{lsa::cosine()}'' directly.

\begin{Shaded}
\begin{Highlighting}[]
\NormalTok{studiedData }\SpecialCharTok{\%\textgreater{}\%} 
  \FunctionTok{as.matrix.data.frame}\NormalTok{() }\SpecialCharTok{\%\textgreater{}\%} \FunctionTok{t}\NormalTok{() }\SpecialCharTok{\%\textgreater{}\%} 
\NormalTok{  lsa}\SpecialCharTok{::}\FunctionTok{cosine}\NormalTok{() }\SpecialCharTok{\%\textgreater{}\%} 
  \FunctionTok{stargazer}\NormalTok{(}\AttributeTok{title =} \StringTok{\textquotesingle{}Cosine Similarity between Cases\textquotesingle{}}\NormalTok{)}
\end{Highlighting}
\end{Shaded}

\begin{table}[!htbp] \centering 
  \caption{Cosine Similarity between Cases} 
  \label{} 
\begin{tabular}{@{\extracolsep{5pt}} cccccccccccc} 
\\[-1.8ex]\hline 
\hline \\[-1.8ex] 
$1$ & $0.985$ & $0.968$ & $0.886$ & $0.626$ & $0.886$ & $0.851$ & $0.632$ & $0.643$ & $0.934$ & $0.970$ & $0.994$ \\ 
$0.985$ & $1$ & $0.994$ & $0.948$ & $0.609$ & $0.901$ & $0.846$ & $0.616$ & $0.597$ & $0.931$ & $0.995$ & $0.997$ \\ 
$0.968$ & $0.994$ & $1$ & $0.974$ & $0.664$ & $0.934$ & $0.881$ & $0.671$ & $0.639$ & $0.950$ & $0.990$ & $0.986$ \\ 
$0.886$ & $0.948$ & $0.974$ & $1$ & $0.662$ & $0.927$ & $0.860$ & $0.669$ & $0.600$ & $0.913$ & $0.954$ & $0.925$ \\ 
$0.626$ & $0.609$ & $0.664$ & $0.662$ & $1$ & $0.876$ & $0.929$ & $1.000$ & $0.984$ & $0.847$ & $0.571$ & $0.609$ \\ 
$0.886$ & $0.901$ & $0.934$ & $0.927$ & $0.876$ & $1$ & $0.989$ & $0.882$ & $0.852$ & $0.991$ & $0.890$ & $0.894$ \\ 
$0.851$ & $0.846$ & $0.881$ & $0.860$ & $0.929$ & $0.989$ & $1$ & $0.934$ & $0.920$ & $0.982$ & $0.828$ & $0.847$ \\ 
$0.632$ & $0.616$ & $0.671$ & $0.669$ & $1.000$ & $0.882$ & $0.934$ & $1$ & $0.984$ & $0.853$ & $0.580$ & $0.616$ \\ 
$0.643$ & $0.597$ & $0.639$ & $0.600$ & $0.984$ & $0.852$ & $0.920$ & $0.984$ & $1$ & $0.843$ & $0.556$ & $0.610$ \\ 
$0.934$ & $0.931$ & $0.950$ & $0.913$ & $0.847$ & $0.991$ & $0.982$ & $0.853$ & $0.843$ & $1$ & $0.916$ & $0.932$ \\ 
$0.970$ & $0.995$ & $0.990$ & $0.954$ & $0.571$ & $0.890$ & $0.828$ & $0.580$ & $0.556$ & $0.916$ & $1$ & $0.990$ \\ 
$0.994$ & $0.997$ & $0.986$ & $0.925$ & $0.609$ & $0.894$ & $0.847$ & $0.616$ & $0.610$ & $0.932$ & $0.990$ & $1$ \\ 
\hline \\[-1.8ex] 
\end{tabular} 
\end{table}

And finally,

\begin{Shaded}
\begin{Highlighting}[]
\FunctionTok{detach}\NormalTok{(currentData)}
\end{Highlighting}
\end{Shaded}


\end{document}
